% !TeX program = lualatex
% !TeX encoding = utf8
% !TeX spellcheck = uk_UA
% !TeX root = ./PyscfBook.tex

\nosection{Список абревіатур}

У цьому посібнику для зручності використовуються стандартні абревіатури з квантової хімії та обчислювальних методів. Наведено їх розшифровку та переклад українською мовою. Абревіатури подано в алфавітному порядку англійських скорочень.

\begin{enumerate}[
    leftmargin=10em,
    font=\bfseries,   % метка жирная
    labelindent=5em,% сдвигает метку вправо от поля
    leftmargin=*,     % использовать всю ширину
]
    \item[HF] Hartree-Fock --- метод Хартрі-Фока.
    \item[RHF] Restricted Hartree-Fock --- обмежений метод Хартрі-Фока (для систем із замкненими оболонками)..
    \item[UHF] Unrestricted Hartree-Fock --- необмежений метод Хартрі-Фока (для систем із відкритими оболонками).
    \item[ROHF] Restricted Open-shell Hartree-Fock --- обмежений метод Хартрі-Фока для відкритих оболонок.
    \item[GHF] Generalized Hartree-Fock --- узагальнений метод Хартрі-Фока.

    \item[SCF] Self-Consistent Field --- самоузгоджене поле.

    \item[CI] Configuration Interaction --- конфігураційна взаємодія.
    \item[CIS] Configuration Interaction Singles --- конфігураційна взаємодія з одноелектронними збудженнями.
    \item[CISD] Configuration Interaction Singles and Doubles --- конфігураційна взаємодія з одно- та двоелектронними. збудженнями.
    \item[CISDT] Configuration Interaction Singles, Doubles and Triples --- конфігураційна взаємодія з одно-, дво- та триелектронними збудженнями.
    \item[FCI] Full Configuration Interaction --- повна конфігураційна взаємодія.
    \item[CASCI] Complete Active Space Configuration Interaction --- конфігураційна взаємодія з повним активним простором.

    \item[MCSCF] Multi-Configurational Self-Consistent Field --- багатоконфігураційне самоузгоджене поле.
    \item[CASSCF] Complete Active Space Self-Consistent Field --- самоузгоджене поле з повним активним простором.
    \item[RASSCF] Restricted Active Space Self-Consistent Field --- самоузгоджене поле з обмеженим активним простором.
    \item[GASSCF] Generalized Active Space Self-Consistent Field --- самоузгоджене поле з узагальненим активним простором.

    \item[CC] Coupled Cluster --- зв'язані кластери.
    \item[CCD] Coupled Cluster Doubles --- метод зв'язаних кластерів з подвійними збудженнями.
    \item[CCSD] Coupled Cluster Singles and Doubles --- метод зв'язаних кластерів з одинарними та подвійними збудженнями.
    \item[CCSD(T)] CCSD with perturbative Triples --- CCSD з пертурбативними потрійними збудженнями.
    \item[CCSDT] Coupled Cluster Singles, Doubles and Triples --- метод зв'язаних кластерів з одинарними, подвійними та потрійними збудженнями.
    \item[CCSDTQ] Coupled Cluster Singles, Doubles, Triples and Quadruples --- метод зв'язаних кластерів з одинарними, подвійними, потрійними та четвірними збудженнями.

    \item[MP] Møller-Plesset perturbation theory --- теорія збурень Мьоллера-Плессета.
    \item[MP2] Second-order Møller-Plesset perturbation theory --- теорія збурень Мьоллера-Плессета другого порядку.
    \item[MP3] Third-order Møller-Plesset perturbation theory --- теорія збурень Мьоллера-Плессета третього порядку.
    \item[MP4] Fourth-order Møller-Plesset perturbation theory --- теорія збурень Мьоллера-Плессета четвертого порядку.
    \item[MP5] Fifth-order Møller-Plesset perturbation theory --- теорія збурень Мьоллера-Плессета п'ятого порядку.

    \item[DFT] Density Functional Theory --- теорія функціонала густини.
    \item[KS-DFT] Kohn-Sham Density Functional Theory --- теорія функціонала густини Кона-Шема.
    \item[TDDFT] Time-Dependent Density Functional Theory --- часозалежна теорія функціонала густини.

    \item[LDA] Local Density Approximation --- наближення локальної густини.
    \item[LSDA] Local Spin Density Approximation --- наближення локальної спінової густини.
    \item[GGA] Generalized Gradient Approximation --- узагальнене градієнтне наближення.
    \item[meta-GGA] meta-Generalized Gradient Approximation --- мета-узагальнене градієнтне наближення.
    \item[hybrid-GGA] hybrid Generalized Gradient Approximation --- гібридне узагальнене градієнтне наближення.
    \item[hybrid-meta-GGA] hybrid meta-Generalized Gradient Approximation --- гібридне мета-узагальнене градієнтне наближення.

    \item[XC] Exchange-Correlation --- обмінно-кореляційний.
    \item[LYP] Lee-Yang-Parr --- кореляційний функціонал Лі-Янга-Парра.
    \item[PW91] Perdew-Wang 91 --- функціонал Пердью-Вонга 1991 року.
    \item[PBE] Perdew-Burke-Ernzerhof --- функціонал Пердью-Берка-Ернцерхофа.
    \item[PBE0] PBE hybrid functional --- гібридний функціонал на основі PBE.
    \item[B3LYP] Becke 3-parameter Lee-Yang-Parr --- гібридний функціонал Бекке з трьома параметрами та кореляційним функціоналом Лі-Янга-Парра.
    \item[B3PW91] Becke 3-parameter Perdew-Wang 91 --- гібридний функціонал Бекке з трьома параметрами та PW91.
    \item[BLYP] Becke-Lee-Yang-Parr --- функціонал з обмінним функціоналом Бекке та кореляційним LYP.
    \item[BP86] Becke-Perdew 86 --- функціонал з обмінним функціоналом Бекке та кореляційним Пердью 86.
    \item[M06] Minnesota 06 functional --- Міннесотський функціонал 2006 року.
    \item[M06-2X] Minnesota 06 double-hybrid functional --- Міннесотський подвійно-гібридний функціонал 2006 року.
    \item[wB97X-D] range-separated hybrid functional with dispersion --- далекосяжний гібридний функціонал з дисперсійною корекцією.
    \item[CAM-B3LYP] Coulomb-Attenuating Method B3LYP --- B3LYP з кулонівським загасанням.

    \item[MRCI] Multi-Reference Configuration Interaction --- багатодетермінантна конфігураційна взаємодія.
    \item[MRCC] Multi-Reference Coupled Cluster --- багатодетермінантний метод зв'язаних кластерів.
    \item[MRPT] Multi-Reference Perturbation Theory --- багатодетермінантна теорія збурень.

    \item[CASPT2] Complete Active Space second-order Perturbation Theory --- теорія збурень другого порядку з повним активним простором.
    \item[RASPT2] Restricted Active Space second-order Perturbation Theory --- теорія збурень другого порядку з обмеженим активним простором.
    \item[NEVPT2] N-Electron Valence state Perturbation Theory --- теорія збурень валентного стану N-електронів другого порядку.

    \item[QMC] Quantum Monte Carlo --- квантовий метод Монте-Карло.
    \item[VMC] Variational Monte Carlo --- варіаційний метод Монте-Карло.
    \item[DMC] Diffusion Monte Carlo --- дифузійний метод Монте-Карло.

    \item[DMRG] Density Matrix Renormalization Group --- метод ренормалізаційної групи матриці густини.

    \item[EOM-CC] Equation-of-Motion Coupled Cluster --- метод зв'язаних кластерів рівнянь руху.
    \item[EOM-CCSD] Equation-of-Motion CCSD --- метод рівнянь руху CCSD.
    \item[SAC-CI] Symmetry-Adapted Cluster Configuration Interaction --- конфігураційна взаємодія кластерів, адаптована до симетрії.

    \item[AO] Atomic Orbital --- атомна орбіталь.
    \item[MO] Molecular Orbital --- молекулярна орбіталь.
    \item[LCAO] Linear Combination of Atomic Orbitals --- лінійна комбінація атомних орбіталей.
    \item[HOMO] Highest Occupied Molecular Orbital --- найвища зайнята молекулярна орбіталь.
    \item[LUMO] Lowest Unoccupied Molecular Orbital --- найнижча незайнята молекулярна орбіталь.
    \item[SOMO] Singly Occupied Molecular Orbital --- однократно зайнята молекулярна орбіталь.

    \item[STO] Slater-Type Orbital --- орбіталь типу Слейтера.
    \item[GTO] Gaussian-Type Orbital --- орбіталь гаусового типу.
    \item[CGTO] Contracted Gaussian-Type Orbital --- контрактована орбіталь гаусового типу.
    \item[PGTO] Primitive Gaussian-Type Orbital --- примітивна орбіталь гаусового типу.

    \item[RI] Resolution of Identity --- розкладання одиниці.
    \item[DF] Density Fitting --- підгонка густини.
    \item[RI-MP2] Resolution of Identity MP2 --- MP2 з розкладанням одиниці.

    \item[BSSE] Basis Set Superposition Error --- похибка суперпозиції базисного набору.
    \item[CP] Counterpoise correction --- контрвагова корекція.

    \item[ZPE] Zero-Point Energy --- енергія нульових коливань.
    \item[ZPVE] Zero-Point Vibrational Energy --- енергія нульових коливань.

    \item[IRC] Intrinsic Reaction Coordinate --- внутрішня координата реакції.
    \item[TS] Transition State --- перехідний стан.
    \item[MEP] Minimum Energy Path --- шлях мінімальної енергії.

    \item[NBO] Natural Bond Orbital --- натуральна зв'язуюча орбіталь.
    \item[NAO] Natural Atomic Orbital --- натуральна атомна орбіталь.
    \item[NPA] Natural Population Analysis --- натуральний популяційний аналіз.

    \item[ESP] Electrostatic Potential --- електростатичний потенціал.
    \item[RESP] Restrained Electrostatic Potential --- обмежений електростатичний потенціал.

    \item[COSMO] COnductor-like Screening MOdel --- модель екранування, подібна до провідника.
    \item[PCM] Polarizable Continuum Model --- модель поляризованого континууму.
    \item[SMD] Solvation Model based on Density --- модель сольватації на основі густини.

    \item[QM/MM] Quantum Mechanics/Molecular Mechanics --- квантова механіка/молекулярна механіка.
    \item[ONIOM] Our own N-layered Integrated molecular Orbital and Molecular mechanics --- багатошаровий інтегрований метод молекулярних орбіталей та молекулярної механіки.

    \item[PBC] Periodic Boundary Conditions --- періодичні граничні умови.
    \item[PW] Plane Wave --- плоска хвиля.

    \item[DFT+U] DFT with Hubbard U correction --- DFT з корекцією Хаббарда U.
    \item[GW] Green's function and screened Coulomb interaction --- функція Гріна та екранована кулонівська взаємодія.
    \item[BSE] Bethe-Salpeter Equation --- рівняння Бете-Солпітера.

    \item[AIMD] Ab Initio Molecular Dynamics --- молекулярна динаміка з перших принципів.
    \item[BOMD] Born-Oppenheimer Molecular Dynamics --- молекулярна динаміка Борна-Оппенгеймера.
    \item[CPMD] Car-Parrinello Molecular Dynamics --- молекулярна динаміка Кара-Паррінелло.

    \item[VCI] Vibrational Configuration Interaction --- конфігураційна взаємодія коливань.
    \item[VPT2] Vibrational second-order Perturbation Theory --- теорія збурень другого порядку для коливань.
    \item[VSCF] Vibrational Self-Consistent Field --- самоузгоджене поле для коливань.

    % Додані абревіатури з посібника (з HFTheory.tex, HFMethod.tex та pyscfBook.pdf)
    \item[post-HF] Post-Hartree-Fock --- пост-методи Хартрі-Фока.
    \item[ERI] Electron Repulsion Integrals --- інтеграли відштовхування електронів.
    \item[J] Coulomb operator --- кулонівський оператор.
    \item[K] Exchange operator --- обмінний оператор.
    \item[F] Fock operator --- оператор Фока.
    \item[VNN] Nuclear-Nuclear Repulsion --- між'ядерне відштовхування.
    \item[D] Density Matrix --- матриця густини.
    \item[DZ] Double-Zeta --- подвійний дзета.
    \item[TZ] Triple-Zeta --- потрійний дзета.
    \item[VQZ] Quadruple-Zeta --- четвірний дзета.
    \item[5Z] Quintuple-Zeta --- п'ятірний дзета.
    \item[IE] Ionization Energy --- енергія іонізації.
    \item[DIIS] Direct Inversion in the Iterative Subspace --- пряма інверсія в ітеративному підпросторі.
\end{enumerate}